\documentclass[runningheads]{llncs}
\usepackage[T1]{fontenc}
\usepackage{booktabs}
\usepackage{array}
\newcolumntype{P}[1]{>{\raggedright\arraybackslash}p{#1}}

\begin{document}
\title{Marchmello01: Topology-Safe Adaptive Lesion Fusion for Interactive PET/CT Segmentation}
\author{Marchmello01}
\authorrunning{Marchmello01}
\institute{Independent researcher}
\maketitle

\begin{abstract}
We describe an interactive whole-body tumor segmentation system for FDG and
PSMA PET/CT. A five-fold nnU-Net residual-encoder ensemble produces the initial
mask, followed by tracer-specific calibration and a prompt-conditioned
distance-transform model. We introduce conservative interaction changes:
clicks are mapped to cropped model space using one crop-origin subtraction;
donor components are validated
independently to prevent a newly detected lesion from hiding a harmful bridge;
PSMA foreground corrections are enabled up to a fixed 128-component burden;
and explicitly annotated PSMA background voxels are removed only when removal
cannot fragment a predicted component. Exact foreground-scribble voxels are
also enforced without dilation when they do not merge accepted lesions. Paired
causal replay on local interactive fixtures increased mean AUC-Dice from
0.8366 to 0.8417 and mean AUC-DMM/F1 from 0.9212 to 0.9216, without an observed
per-case regression. The method
is stateless, deterministic with respect to cumulative scribbles, and is
released with complete upstream attribution.
\keywords{autoPET V \and interactive segmentation \and PET/CT \and topology}
\end{abstract}

\section{Introduction}
autoPET V evaluates lesion segmentation across FDG and PSMA PET/CT under
iterative simulated and clinician-driven scribbles. High voxel overlap alone
is insufficient: a useful method must also create missed lesion instances and
remove false-positive instances without merging anatomically separate tumors.
Our submission targets this topological failure mode while preserving strong
initial segmentation from public AutoPET-III and AutoPET-IV components.

\section{Methods}
\subsection{Data}
\subsubsection{Training data}
The pretrained components use the public data and preprocessing described by
the upstream AutoPET-III LesionTracer and AutoPET-IV LesionLocator releases.
Our changes introduce no additional training data or weight optimization.

\subsubsection{Validation data}
Local validation used four FDG and four PSMA interactive fixtures, with
cumulative challenge-format scribbles replayed for five steps. These fixtures
were used for paired engineering comparisons and are not presented as an
unbiased estimate of hidden-test rank.

\subsection{Data preprocessing}
CT and PET are resampled and normalized by the frozen nnU-Net plans. Predictions
are returned to the CT reference grid, and the output copies its size, spacing,
origin, and direction. Input-grid click coordinates are mapped into cropped
model space by subtracting the crop origin once and then applying the
crop-to-resampled scale. No metadata identifying the tracer is required.

\subsection{Initial segmentation and tracer routing}
Three frozen classifiers route each scan to FDG or PSMA. The initial mask is a
five-fold AutoPET-III LesionTracer ResEncL ensemble without mirror test-time
augmentation. FDG uses a probability threshold of 0.47 and burden-adaptive
small-component filtering. PSMA uses threshold 0.50, five-voxel filtering,
and the upstream frozen logistic component pruner at threshold 0.86.

\subsection{Algorithm and topology-safe interaction}
Later calls reconstruct the initial mask and run the fold-0 EDT LesionLocator
with all cumulative scribbles. Each clicked foreground donor component is
considered independently. It is accepted only if union with the current mask
does not reduce the connected-component count, preventing lesion bridges.
For PSMA, foreground correction requires at least six cumulative foreground
points and at most 128 initial components. Each connected background stroke is
also handled independently: only its explicitly marked voxels are erased, and
the erase is accepted only when it does not increase component count. FDG
background behavior remains unchanged from the upstream system.
Finally, each connected foreground scribble stroke is treated as direct
supervision: its exact marked voxels are added without dilation if doing so
does not reduce component count. This last operation is independent of network
probability calibration.

\subsection{Data postprocessing}
FDG predictions use a probability threshold of 0.47 and burden-adaptive
connected-component filtering. PSMA uses threshold 0.50, five-voxel filtering,
and a frozen logistic component pruner with false-positive threshold 0.86.
All interactive unions and erasures are checked using 18-connected components.

\subsection{Training and test parameters}
All network weights are unchanged pretrained weights. Consequently, no new
optimizer, schedule, augmentation, or training hardware is applicable to this
derivative. Inference uses five initial-model folds, one interactive fold,
disabled mirror TTA, binary connected components, and cumulative stateless
scribble replay. The submitted image targets an NVIDIA T4 with 16~GB VRAM.

\subsection{GitHub repository}
Code, tests, build recipes, notices, and checkpoint integrity hashes are
released under a permissive license at
github.com/Marchmello01/\allowbreak autopet-v-topology-safe-adaptive.

\section{Results}
Relative to the frozen v2 policy, exact foreground-stroke enforcement raised
mean causal AUC-Dice from 0.8366 to 0.8417 and mean AUC-DMM/F1 from 0.9212 to
0.9216 over seven defined local cases. Dice improved in four cases and was
unchanged in three; F1 improved in one and was unchanged in six. The largest
Dice change was PSMA3 (0.7937 to 0.8238), while PSMA1 F1 changed from 0.8871 to
0.8898. These local fixtures support selection of the policy but are not a
substitute for hidden-test evaluation. The earlier frozen v2 container obtained preliminary-test Dice
0.854649, F1/DMM 0.834179, and mean position 3.5. The preliminary set contains
only five cases and is reported as an implementation check, not as a final-test
performance claim.

\section{Discussion and conclusion}
The method changes interaction logic rather than claiming newly trained model
weights. Its main advantage is a simple invariant: one useful new lesion may
not compensate for a separate harmful bridge in the same update. The
128-component ceiling is deliberately conservative and should be re-estimated
on a larger independent validation set. Exact foreground-stroke enforcement
reduces dependence on probability calibration, but the small validation cohort
does not establish final-set performance.

\begin{table}[ht]
\caption{Algorithm details.}\label{tab1}
\small
\begin{tabular}{P{0.18\textwidth}P{0.18\textwidth}P{0.18\textwidth}P{0.18\textwidth}P{0.18\textwidth}}
\toprule
Team & Algorithm & Preprocessing & Postprocessing & Augmentation \\
\midrule
Marchmello01 & Tracer Adaptive LesionLocator Ensemble & nnU-Net PET/CT normalization and resampling & Tracer calibration; topology-safe component fusion & Upstream pretrained models; no new training \\
\bottomrule
\end{tabular}

\vspace{1em}
\begin{tabular}{P{0.18\textwidth}P{0.18\textwidth}P{0.18\textwidth}P{0.18\textwidth}P{0.18\textwidth}}
\toprule
TTA & Ensembling & Framework & Architecture & Training loss \\
\midrule
None & Five-fold K0; one interactive fold; three-model tracer vote & nnU-Net v2 / PyTorch & 3-D residual-encoder U-Nets & Unchanged upstream Dice and cross-entropy objectives \\
\bottomrule
\end{tabular}

\vspace{1em}
\begin{tabular}{P{0.23\textwidth}P{0.23\textwidth}P{0.23\textwidth}P{0.23\textwidth}}
\toprule
Training data & Dimensionality & Pretrained models & Hardware \\
\midrule
Public AutoPET FDG/PSMA cohorts; no new training & 3-D whole-body PET/CT with sliding-window inference & Public AutoPET-III initializer and AutoPET-IV interaction model & NVIDIA T4 16~GB inference target \\
\bottomrule
\end{tabular}
\end{table}

\begin{credits}
\subsubsection{Acknowledgements}
This work uses the public AutoPET-III LesionTracer and AutoPET-IV LesionLocator
projects. We thank
their authors and the autoPET organizers.
\subsubsection{Disclosure of Interests}
The author declares no competing interests.
\end{credits}

\begin{thebibliography}{3}
\bibitem{nnunet}
Isensee, F., Jaeger, P.F., Kohl, S.A.A., Petersen, J., Maier-Hein, K.H.:
nnU-Net: a self-configuring method for deep learning-based biomedical image
segmentation. Nature Methods \textbf{18}, 203--211 (2021)
\bibitem{autopetsoftware}
MIC-DKFZ: AutoPET-III LesionTracer and AutoPET-IV LesionLocator, Apache-2.0
software releases (2024--2025)
\bibitem{curriculum}
Three-Phase Scribble-Adaptive Curriculum Learning for autoPETV Grand
Challenge. arXiv:2608.22096 (2026)
\end{thebibliography}
\end{document}
